% =====================================================================
%  Thème 5 — Résolution des triangles — NIVEAU FACILE
%  Corrigé
% =====================================================================
\documentclass[11pt,a4paper]{article}
\usepackage[utf8]{inputenc}
\usepackage[T1]{fontenc}
\usepackage{lmodern}
\usepackage[french]{babel}
\usepackage{amsmath,amssymb,mathtools}
\usepackage{icomma}
\usepackage{xcolor}
\usepackage{colortbl}
\usepackage{tikz}
\usepackage[most]{tcolorbox}
\usepackage{enumitem}
\usepackage{array}
\usepackage[a4paper,margin=2.2cm]{geometry}
\usetikzlibrary{arrows.meta,calc}

\definecolor{primary}{RGB}{214,51,108}
\definecolor{primarydark}{RGB}{140,20,64}
\definecolor{excol}{RGB}{0,135,90}

\DeclareMathOperator{\tg}{tg}
\DeclareMathOperator{\cotg}{cotg}
\newcommand{\degr}{^{\circ}}
\newcommand{\Z}{\mathbb{Z}}

\newcounter{exo}
\newcommand{\exo}[1]{\medskip\refstepcounter{exo}%
  \noindent{\large\bfseries\color{primarydark}Exercice \theexo}%
  \ \textcolor{primary}{\textbullet}\ \textbf{#1}\par\nopagebreak\smallskip}
\newcommand{\rep}{\textcolor{excol}{$\blacktriangleright$}\ }

\setlength{\parindent}{0pt}
\pagestyle{empty}

\begin{document}

\begin{center}
{\LARGE\bfseries\color{primary}Thème 5 — Résolution des triangles}\\[2pt]
{\color{primarydark}\textsc{Niveau Facile — Corrigé détaillé}}
\end{center}
\hrule height 1pt
\medskip

\exo{Calculer une aire ($\mathcal{A}=\tfrac12\,ab\sin\hat C$)}
On multiplie les deux côtés donnés par le sinus de l'angle compris, divisé par $2$.
\begin{itemize}[leftmargin=1.4em,itemsep=1pt]
\item[\textbf{(a)}] \rep $\mathcal{A}=\dfrac12\cdot6\cdot8\cdot\sin30\degr=24\cdot\dfrac12=12$.
\item[\textbf{(b)}] \rep $\mathcal{A}=\dfrac12\cdot5\cdot4\cdot\sin90\degr=10\cdot1=10$.
\item[\textbf{(c)}] \rep $\mathcal{A}=\dfrac12\cdot10\cdot6\cdot\sin60\degr=30\cdot\dfrac{\sqrt3}{2}=15\sqrt3$.
\item[\textbf{(d)}] \rep $\mathcal{A}=\dfrac12\cdot8\cdot8\cdot\sin45\degr=32\cdot\dfrac{\sqrt2}{2}=16\sqrt2$.
\item[\textbf{(e)}] \rep $\mathcal{A}=\dfrac12\cdot3\cdot4\cdot\sin120\degr=6\cdot\dfrac{\sqrt3}{2}=3\sqrt3$
  (rappel $\sin120\degr=\sin60\degr=\tfrac{\sqrt3}{2}$).
\end{itemize}

\exo{Trouver le 3\textsuperscript{e} côté (loi des cosinus)}
On applique $c^2=a^2+b^2-2ab\cos\hat C$.
\begin{itemize}[leftmargin=1.4em,itemsep=1pt]
\item[\textbf{(a)}] \rep $c^2=25+64-2\cdot40\cdot\cos60\degr=89-80\cdot\tfrac12=89-40=49$, donc $c=7$.
\item[\textbf{(b)}] \rep $\cos120\degr=-\tfrac12$ : $c^2=9+25-2\cdot15\cdot(-\tfrac12)=34+15=49$, donc $c=7$.
\item[\textbf{(c)}] \rep $\cos90\degr=0$ : $c^2=36+64=100$, donc $c=10$ (c'est Pythagore).
\item[\textbf{(d)}] \rep $c^2=4+9-2\cdot6\cdot\tfrac12=13-6=7$, donc $c=\sqrt7$.
\end{itemize}

\exo{Trouver un côté manquant (loi des sinus)}
On isole le côté cherché : le côté est proportionnel au sinus de l'angle opposé.
\begin{itemize}[leftmargin=1.4em,itemsep=1pt]
\item[\textbf{(a)}] \rep $\dfrac{\sin30\degr}{6}=\dfrac{\sin90\degr}{b}\Rightarrow
   b=6\cdot\dfrac{\sin90\degr}{\sin30\degr}=6\cdot\dfrac{1}{1/2}=12$.
\item[\textbf{(b)}] \rep $\dfrac{\sin30\degr}{5}=\dfrac{\sin45\degr}{c}\Rightarrow
   c=5\cdot\dfrac{\sin45\degr}{\sin30\degr}=5\cdot\dfrac{\sqrt2/2}{1/2}=5\sqrt2$.
\item[\textbf{(c)}] \rep $\dfrac{\sin60\degr}{3}=\dfrac{\sin90\degr}{b}\Rightarrow
   b=3\cdot\dfrac{1}{\sqrt3/2}=\dfrac{6}{\sqrt3}=2\sqrt3$.
\end{itemize}

\exo{Retrouver un angle à partir des trois côtés}
On calcule $\cos\hat C=\dfrac{a^2+b^2-c^2}{2ab}$ et on lit l'angle remarquable.
\begin{itemize}[leftmargin=1.4em,itemsep=1pt]
\item[\textbf{(a)}] \rep $\cos\hat C=\dfrac{9+25-49}{2\cdot15}=\dfrac{-15}{30}=-\dfrac12
   \Rightarrow \hat C=120\degr$.
\item[\textbf{(b)}] \rep $\cos\hat C=\dfrac{25+64-49}{2\cdot40}=\dfrac{40}{80}=\dfrac12
   \Rightarrow \hat C=60\degr$.
\item[\textbf{(c)}] \rep $\cos\hat C=\dfrac{36+64-100}{2\cdot48}=\dfrac{0}{96}=0
   \Rightarrow \hat C=90\degr$ (triangle rectangle).
\item[\textbf{(d)}] \rep $\cos\hat C=\dfrac{4+9-7}{2\cdot6}=\dfrac{6}{12}=\dfrac12
   \Rightarrow \hat C=60\degr$.
\end{itemize}

\exo{Le triangle rectangle (SOH-CAH-TOA et Pythagore)}
\begin{itemize}[leftmargin=1.4em,itemsep=2pt]
\item[\textbf{(a)}] \rep Côté opposé $=\text{hyp}\times\sin60\degr=6\cdot\dfrac{\sqrt3}{2}=3\sqrt3$ ;
   côté adjacent $=\text{hyp}\times\cos60\degr=6\cdot\dfrac12=3$.
\item[\textbf{(b)}] \rep Par Pythagore, hypoténuse $=\sqrt{5^2+5^2}=\sqrt{50}=5\sqrt2$. Le triangle
   étant isocèle, ses deux angles aigus sont égaux et de somme $90\degr$, donc chacun vaut $45\degr$.
\item[\textbf{(c)}] \rep Hypoténuse $=\sqrt{5^2+12^2}=\sqrt{25+144}=\sqrt{169}=13$.
\end{itemize}

\end{document}
